\documentclass[11pt]{article}
\usepackage[margin=1in]{geometry}
\usepackage{amsmath,amssymb,amsthm,hyperref,enumitem}
\hypersetup{colorlinks=true,linkcolor=blue,urlcolor=blue}
\setlength{\emergencystretch}{1em}
\newtheorem{theorem}{Theorem}
\newcommand{\E}{\mathbb E}
\newcommand{\Pp}{\mathbb P}
\newcommand{\Var}{\operatorname{Var}}
\newcommand{\Cov}{\operatorname{Cov}}
\newcommand{\N}{\mathbb N}
\newcommand{\R}{\mathbb R}
\title{All-gap moments and Gaussian limits\\for ballistic-deposition roots}
\author{Principia Math In Progress}
\date{22 September 2026}
\begin{document}
\maketitle
\noindent\textbf{MathDB 342282. Full intended theorem.}
The complete accepted manuscript and current preprint have been read;
the refreshed bounded source and claim audit is recorded with this proof.

\section*{Task 1: statement and scope}
Consider the cyclic strip with \(K\) sites. At each step a uniform site
\(j\) is selected and its height becomes
\(1+\max(H(j-1),H(j),H(j+1))\); other heights are unchanged.
Initially all heights are zero. A root is a site that ever receives a
particle at height one. For \(i\geq1\), \(D_{i,K}\) counts consecutive
roots at cyclic distance \(g=i+1\), including the wrapping pair.

The first conjecture in Section 3.5 of Mansour--Rastegar--Roitershtein
[1] asks for linear moments and a normal limit. It prints
\(\sigma(D_{i,K})=b_iK\), although \(\sigma\) denotes standard deviation.
The surrounding discussion and Theorems 3.3 and 3.5 concern linear
\emph{variance}. We prove that intended all-\(i\) assertion and both
Gaussian conjectures. The elementary variance/standard-deviation
wording discrepancy earns no research credit.

\begin{theorem}
For every fixed \(i\geq1\), there are positive rational numbers \(a_i,b_i\)
such that, whenever \(K\geq2(i+4)\),
\[
 \E D_{i,K}=a_iK,\qquad \Var(D_{i,K})=b_iK.                 \tag{1}
\]
Moreover,
\[
 \frac{D_{i,K}-a_iK}{\sqrt{b_iK}}\ \Longrightarrow\ N(0,1).
                                                               \tag{2}
\]
The vector of all coordinates standardized by their actual means and
standard deviations converges weakly in
\(\R^{\N}\), with its product topology, to a centered Gaussian process.
For the finitely small sizes where a coordinate has zero variance,
define that standardized coordinate as zero; also set coordinates
\(i\geq K\) to zero.
\end{theorem}

\section*{Independent priorities and local gap indicators}
Run independent rate-one Poisson clocks at the sites. Their merged
sequence of site labels is independent and uniform, so it gives the
same deposition process. Its first arrival times are independent
exponentials. A root forms at \(j\) exactly when its first arrival
precedes the first arrivals at both neighbors: otherwise a neighbor
already has positive height. Later arrivals cannot create a root at
a previously occupied site. Applying the exponential distribution
function gives independent uniform priorities \(U_j\) and the exact
joint representation
\[
 \xi_j={\bf1}\{U_j<U_{j-1},\ U_j<U_{j+1}\}.
\]
All indices here are cyclic. If \(K>g+2\), put
\[
 Y_j^{(g)}=\xi_j\xi_{j+g}
              \prod_{t=1}^{g-1}(1-\xi_{j+t}).
 \qquad D_{i,K}=\sum_{j=1}^K Y_j^{(g)}.                       \tag{3}
\]
The indicator depends only on
\(U_{j-1},\ldots,U_{j+g+1}\), a window of length
\(m=g+3=i+4\).

Use the same definition on a two-sided infinite independent sequence.
Its indicators are stationary and \((m-1)\)-dependent: disjoint
priority windows give independent families. Define
\[
 a_i=\E Y_0^{(g)},\qquad
 b_i=\Var(Y_0^{(g)})
       +2\sum_{h=1}^{m-1}\Cov(Y_0^{(g)},Y_h^{(g)}).           \tag{4}
\]
Each probability in (4) is a finite sum of equally likely strict
orders of finitely many continuous priorities. Thus \(a_i,b_i\) are
rational. A configuration with low values at \(0,g\), high external
neighbors, and increasing interior values has exactly those two
roots in \([0,g]\). It occupies a nonempty open set, so \(a_i>0\).

For \(K\geq2m\), a single window has the line law. A pair of overlapping
windows has the line law at offset \(h\), or at offset \(-h\) across the
cyclic seam. All other pairs are independent. Summing covariances gives
(1) with the constants (4), except that positivity of \(b_i\) still
requires proof.

\section*{Strictly positive variance density}
Fix \(g\geq2\), and write \(q=2g-1\), \(L=2g+3\).
Partition \(\lfloor K/L\rfloor\) disjoint consecutive blocks of length
\(L\), leaving any remaining sites aside. Label sites within each block
\(-1,0,1,\ldots,2g,2g+1\). Its anchors are \(0,2g\), guards are
\(-1,2g+1\), and interior sites are \(1,\ldots,q\).
A block is good if both anchors have priority below \(1/4\), both guards
have priority above \(3/4\), and all interior priorities exceed \(1/4\).
Its probability is
\[
 p=4^{-4}(3/4)^q>0.                                         \tag{5}
\]

Reveal all anchor, guard, and leftover priorities, and for every block
reveal whether all its interior priorities exceed \(1/4\). In blocks
that are not good, reveal the whole interior as well. Conditional on
this information, the interiors of the good blocks are independent;
in each good block they are independent uniforms on \((1/4,1)\).
Both anchors are forced roots and both guards are forced nonroots.
All roots outside the good interiors are determined.

Consequently \(D_{i,K}\), conditionally, is a constant plus the sum
of independent quantities \(H\), one per good block, where \(H\)
counts gaps of length \(g\) inside its anchored segment \([0,2g]\).
The fixed anchor roots separate all such gaps from exterior gaps.
Two interior order patterns have different values of \(H\).
If \(U_1<\cdots<U_q\), there are no interior roots, so \(H=0\).
For the second pattern make \(U_g\) the least interior priority,
make \(U_1,\ldots,U_{g-1}\) increase, and make
\(U_{g+1},\ldots,U_q\) decrease. Choose one fixed total order with these
properties. Its only interior root is \(g\), so \(H=2\).
These descriptions also hold when \(g=2\).

Each of the two selected total orders has conditional probability
\(1/q!\). For an independent conditional copy \(H'\),
\[
 \Var(H\mid\text{revealed})
 =\tfrac12\E((H-H')^2\mid\text{revealed})
 \ \geq\frac4{(q!)^2}.                                     \tag{6}
\]
Conditional independence, total variance, and (5) now give
\[
 \Var(D_{i,K})\geq
 \left\lfloor\frac K L\right\rfloor
                  \frac{4p}{(q!)^2}.
\]
Combining this with (1) and sending \(K\) to infinity proves the explicit
positive bound
\[
 b_i\geq \frac{4^{\, -3}(3/4)^{\,2g-1}}
                    {(2g+3)((2g-1)!)^2}>0.                 \tag{7}
\]

\section*{The scalar central limit theorem}
We include a blocking argument. On the infinite line let
\(X_j=Y_j^{(g)}-a_i\) and \(S_K=\sum_{j=1}^K X_j\).
Finite dependence gives
\(\Var(S_K)=Kb_i+O(1)\). Take consecutive long blocks of length
\(\ell=\lfloor K^{1/4}\rfloor\), separated by \(m-1\) discarded
indices, and discard the final incomplete block. The long-block sums
\(B_r\) are independent and centered, and \(|B_r|\leq\ell\).
The discarded sum \(R_K\) has variance
\[
 \Var(R_K)=O(K/\ell+\ell)=o(K),
\]
since each discarded index has at most \(2m-2\) dependent neighbors.
The \(L^2\) triangle inequality yields
\(s_K^2:=\sum_r\Var(B_r)\sim b_iK\).

For completeness, the independent-block normal limit follows from
the Taylor expansion of characteristic functions. For each fixed real
\(t\), the total third-order remainder is bounded by
\[
 \frac{|t|^3}{6s_K^3}\sum_r\E|B_r|^3
 \leq\frac{|t|^3\ell}{6s_K}\longrightarrow0.
\]
Also \(\max_r\Var(B_r)/s_K^2\leq\ell^2/s_K^2\to0\).
Multiplying the characteristic functions, or taking their logarithms
near one, therefore gives
\(\E\exp(it\sum_r B_r/s_K)\to e^{-t^2/2}\).
Thus \(S_K/\sqrt{b_iK}\) converges to \(N(0,1)\).

Couple the cycle to the line by using the same priorities at sites
\(1,\ldots,K\). Only a bounded number, depending on \(m\), of indicator
windows cross the seam. The cyclic and line sums therefore differ by
a deterministically bounded amount. Dividing by \(\sqrt{b_iK}\)
proves (2).

\section*{The joint product-space limit}
For fixed \(r\), apply the same argument to any real linear combination
of the first \(r\) local indicators, centered and divided by
\(\sqrt{b_i}\). This remains bounded and finitely dependent.
Its variance density is nonnegative. If positive, the blocking proof
applies after rescaling; if zero, its line-sum variance is \(O(1)\),
so its sum divided by \(\sqrt K\) converges to zero in \(L^2\).
The bounded seam correction still vanishes.

The Cramer--Wold criterion gives a centered Gaussian limit for every
finite set of coordinates, with covariance matrix
\[
 \Sigma_{ij}=\frac{1}{\sqrt{b_i b_j}}
       \sum_{h\in\mathbb Z}
       \Cov\bigl(Y_0^{(i+1)},Y_h^{(j+1)}\bigr).              \tag{8}
\]
Each sum is finite. These matrices are consistent and positive
semidefinite, as limits of covariance matrices; their diagonals are one.
They define a Gaussian law on the countable product.

Every standardized coordinate has second moment at most one.
Given \(\varepsilon>0\), choose \(M_i>0\) with
\(\sum_i M_i^{-2}<\varepsilon\). By Chebyshev and the union bound,
all these vector laws give mass at least \(1-\varepsilon\) to the
compact product \(\prod_i[-M_i,M_i]\). They are uniformly tight.
Finite-dimensional convergence and tightness prove convergence in
the product topology to the Gaussian law with covariance (8).

\section*{Tasks 2--5: four proposed mathematical follow-ups}
\begin{enumerate}[label=\textbf{Task \arabic*:},start=2,leftmargin=*]
\item Derive closed formulas, or generating functions, for all
mean and variance constants \(a_i,b_i\) in (4).
\item Prove a quantitative normal approximation and a local central
limit theorem for \(D_{i,K}\), with explicit dependence on fixed \(i\).
\item Classify Gaussian and Poisson limiting regimes when the gap
index \(i=i(K)\) grows with the strip size.
\item Determine sufficient conditions on nonuniform site arrival
rates for analogous linear-variance and joint Gaussian limits.
\end{enumerate}
These are proposed directions, not additional claimed solutions.

\noindent\textbf{Recorded verification.}
Four exact groups check 1,093 adjacent-order patterns; evaluate
the constants for \(1\leq i\leq10\), agreeing with all seven known
pairs; check 409,112 first-arrival orders modulo rotation for
\(3\leq K\leq10\); and check variance witnesses for \(2\leq g\leq50\),
including exhaustive conditional variances for \(g=2,3,4\).
These finite computations supplement the proof; they do not establish
an infinite-parameter statement or a normal limit by numerical fitting.
The proof has been checked only by its authoring agent, without
independent peer review or formal proof-assistant verification.

\section*{Sources and status}
[1] T. Mansour, R. Rastegar, A. Roitershtein.\\
\emph{On ballistic deposition process on a strip}.\\
\href{https://arxiv.org/abs/1903.12548}{arXiv:1903.12548v2}
(19 June 2019), all 27 PDF pages read.
The target is Section 3.5, p.~23; the gap definition is on p.~10.
\par\smallskip\noindent Journal article: \emph{J. Stat. Phys.} 177 (2019), 626--650,
\href{https://doi.org/10.1007/s10955-019-02383-4}
{journal DOI}. The complete
\href{https://pmc.ncbi.nlm.nih.gov/articles/PMC6913540/}{accepted manuscript,
PMC6913540 / NIHMS1058046}, including Sections 1--3.5, all three tables,
references, and all seven figures, was read on 22 September 2026.
Its Section 3.5 retains both target conjectures.
\href{https://pmc.ncbi.nlm.nih.gov/about/authorms/}{PMC's version policy}
identifies this as the manuscript accepted after peer review.
The separately typeset publisher version remains unread; no identity
with that version is asserted. The accepted final supplies the
source-scope check, replacing the earlier pending assessment.

\medskip\noindent
Status at 22 September 2026: complete intended mathematical result;
eligible within the recorded bounded audit. Initial and refreshed MathDB and
bounded public-claim checks are recorded in the package. No credit is
claimed for correcting the printed standard-deviation formula, for
known finite-\(i\) moment calculations, or separately for the joint
strengthening.
\end{document}
